
\section{A correction reaches different decision makers at different times}
\label{sim:actor-vintage-decision}
Replaying the corrected history answers an inference question. It does not
yet say which decision maker could use the correction. Suppose the analyst
learns on 1 February 2026 that the earlier observation was wrong, while the
manager receives that correction on 1 March. On 15 February they possess
different information even though they are studying the same company. A
simulation of the manager's choice must use the manager's information.

Continue the hypothetical fixed-versus-floating example. Under the original
observation the posterior is $(1,0)$; under the correction it is
$(1/4,3/4)$. Suppose an unhedged exposure costs zero for the fixed type and
20 for the floating type. A perfect hedge costs 6 in either case. These are
supplied losses in the same units, and both actions are assumed feasible.
They are enough to define this comparison without inventing an observed
hedging price or an accounting transaction.

Let $p$ denote the probability of the floating type under the particular
actor's available information. The expected losses are therefore
\begin{equation}
  L(\text{unhedged}\mid p)=20p,\qquad
  L(\text{hedged}\mid p)=6.
  \label{sim:actor-hedge-loss}
\end{equation}
The original posterior gives an unhedged loss of zero. The corrected
posterior gives 15, so the hedge has the smaller declared expected loss.
At $p=3/10$ the two losses both equal 6 and both actions remain minimizers.
An implementation must retain that tie unless a separate choice rule is
supplied. The manager on 15 February still uses the old posterior, while
the analyst can already select the hedge under the corrected one.

The selected information procedure takes an actor and an exact cutoff,
selects the available unreplaced observation in each information slot, and
passes that request to the finite evaluator. Adding the future correction
does not change the manager's stored 15 January request or its decision.
Later knowledge can motivate a new decision, but it must not silently
replace the information used to reproduce an earlier one.

Unknown availability cannot be treated as known absence. Two incomparable
available revisions also leave the request unresolved; choosing the newer
file arbitrarily would add an authority rule. The checked examples retain
these distinctions, refuse corrupt observation content and fabricated
posterior results, and leave unknown action losses unresolved. They select
actions from the declared feasible set; they do not execute a hedge or infer
its causal effect. Actual actor access, registry completeness, instrument
terms and feasible financing still need evidence.
